\begin{summarybox}
	\textbf{\textcolor{CSummary}{一句话核心}}

	\medskip
	判别式法通过构造二次方程，将值域问题转化为含参二次方程有实根的条件，利用判别式非负解出 $y$ 的范围。

	\medskip
	\textbf{\textcolor{CSummary}{使用条件}}
	\begin{itemize}[leftmargin=*,itemsep=0.5em,topsep=0.4em]
		\item \textbf{条件 1（分子分母二次）：} 函数为 $\displaystyle y=\frac{ax^2+bx+c}{dx^2+ex+f}$，且 $d\neq0$。
		\item \textbf{条件 2（定义域处理）：} 若分母 $\Delta_d<0$，定义域为 $\mathbb{R}$，可直接使用；若 $\Delta_d\ge0$，需先确定定义域，再结合判别式法。
	\end{itemize}

	\medskip
	\textbf{\textcolor{CSummary}{关键提醒}}
	\begin{itemize}[leftmargin=*,itemsep=0.5em,topsep=0.4em]
		\item \textbf{易错点（系数退化）：} 当 $a-dy=0$ 时务必单独检验 $y$ 是否可取，不可遗漏。
		\item \textbf{检查项（符号展开）：} 展开判别式时注意多项式乘法符号，建议列出中间项。
	\end{itemize}

\end{summarybox}
